\documentclass[11pt]{amsart}
\usepackage[T1]{fontenc}
\usepackage[utf8]{inputenc}
\usepackage{lmodern}
\usepackage[margin=24mm]{geometry}
\usepackage{microtype}
\usepackage{amsmath,amssymb,amsthm,mathtools,booktabs,hyperref}
\hypersetup{colorlinks=true,linkcolor=blue,urlcolor=blue}
\newtheorem{theorem}{Theorem}
\newtheorem{lemma}{Lemma}
\newtheorem{corollary}{Corollary}
\newtheorem{proposition}{Proposition}
\theoremstyle{remark}\newtheorem{remark}{Remark}
\newtheorem{problem}{Problem}
\newcommand{\N}{\mathbb N}
\newcommand{\Syr}{\operatorname{Syr}}
\newcommand{\Pass}{\operatorname{Pass}}
\newcommand{\odd}{\operatorname{odd}}
\newcommand{\Psh}{\mathbb P_{\mathrm{sh}}}
\newcommand{\Ph}{\mathbb P_{\mathrm{Haar}}}
\title[Local first-passage progressions]{An unconditional local first-passage progression estimate\\for Syracuse entrances and short upward excursions}
\author{Research working note}
\date{18 November 2026 -- preliminary note, not a proof of the Collatz conjecture}
\begin{document}
\begin{abstract}
We prove a local, unconditional power saving for the probability of an upward shortcut-Collatz excursion within $\lfloor\log X\rfloor$ steps \emph{after} a first Syracuse passage below $B=\lfloor X^{99/100}\rfloor$, when the original odd starting integer is sampled uniformly from $(B,2B]$.  Exact valuation-word cylinders map consecutive source points to long arithmetic progressions of entry states; for typical short first passages this implies $2$-adic residue balance.  Chernoff estimates handle the long and high-valuation first passages.  Explicit parameters give an $O(X^{-1/200})$ exceptional probability.  The result also applies to harmonic weights on this source shell.  Importantly, it is a \emph{local-shell} result, not a bound for the limiting logarithmic first-passage law in the source manuscript, and it does not control arbitrarily late returns.
\end{abstract}
\maketitle

\section{Definitions and the elementary progression mechanism}
For positive integers put
\[
 T(n)=\begin{cases}n/2 &2\mid n,\\(3n+1)/2&2\nmid n,\end{cases}
 \qquad \Syr(n)=\frac{3n+1}{2^{\nu_2(3n+1)}}\quad(n\text{ odd}).
\]
Fix $X>1$, $0<\beta<1/2$, and abbreviate
\[
 b=1-2\beta,\qquad B=\lfloor X^b\rfloor,\quad Y=X^{1-\beta},\quad
 L=\lfloor\log X\rfloor,\quad m=2^{L-1}.
\]
Our \emph{local shell} is $\mathcal I_B=(B,2B]\cap(2\N+1)$, with $R=\#\mathcal I_B=B/2+O(1)$, equipped with normalized uniform counting measure $\Psh$.  Put
\[
 \tau_B(n)=\inf\{k\ge1:\Syr^k(n)\le B\},\qquad
 \Pass_B(n)=\Syr^{\tau_B(n)}(n)
\]
when $\tau_B(n)<\infty$.  The full limiting entrance law obtained from starts tending to infinity is \emph{not} the same as this local law.

\begin{lemma}[Valuation cylinder and interval progression]\label{lem:progression}
Let $\mathbf a=(a_1,\ldots,a_k)$ be a positive-integer valuation word and $A=a_1+\cdots+a_k$.  Among odd positive integers, those whose first $k$ Syracuse valuations equal $\mathbf a$ form exactly one residue class modulo $2^{A+1}$.  On that class, consecutive inputs $n=n_0+2^{A+1}j$ have
\[
 \Syr^i(n)=\Syr^i(n_0)+2\,3^i\,2^{A-A_i}j\quad(0\le i\le k),
 \qquad A_i=\sum_{\ell\le i}a_\ell,
\]
with $A_0=0$ (and for $i=0$ the displayed increment is $2^{A+1}j$).  In particular, the entry outputs $\Syr^k(n)$ form an arithmetic progression of step $2\cdot3^k$.  The additional restrictions
\[
 n\in\mathcal I_B,\qquad \Syr^i(n)>B\ (0\le i<k),\qquad
 \Syr^k(n)\le B
\]
select a \emph{consecutive interval} of indices $j$.  Consequently, for any odd residue $r\bmod2^L$ and any such valuation word $\mathbf a$,
\[
 \left|\#\{\text{admissible inputs giving output }r\}-\frac{N_{\mathbf a}}m\right|\le1,
\]
where $N_{\mathbf a}$ is the total number of admissible inputs on that cylinder.
\end{lemma}
\begin{proof}
The Syracuse valuation coding is bijective between words with total valuation $A$ and their compatible odd residue modulo $2^{A+1}$; this follows inductively by imposing $\nu_2(3\Syr^{i-1}(n)+1)=a_i$. On a fixed word, each iterate is an affine function of $n$ with positive slope $3^i/2^{A_i}$.  Translating $n$ by $2^{A+1}$ therefore translates the $i$-th iterate by $3^i2^{A+1-A_i}$, yielding the formula above.  Every first-passage inequality is a half-line in $j$, so their intersection with the input interval is consecutive.  As $3^k$ is invertible modulo $m=2^{L-1}$, the final outputs $y_0+2\cdot3^kj$ run through all odd residues modulo $2^L$ once every $m$ consecutive $j$'s.  Each count differs from $N_{\mathbf a}/m$ by less than one.
\end{proof}
\begin{lemma}[Uniform residue balance for truncated passage]\label{lem:balance}
Fix positive integers $H,A_0$.  Define $G$ as the starts $n\in\mathcal I_B$ satisfying $\tau_B(n)\le H$ and $\sum_{i\le\tau_B(n)}\nu_2(3\Syr^{i-1}(n)+1)\le A_0$.  Let $q_G(r)=R^{-1}\#\{n\in G:\Pass_B(n)\equiv r\pmod{2^L}\}$, $a_G=\Psh(G)$. Then
\begin{equation}\label{eq:residue}
 \boxed{\left|q_G(r)-\frac{a_G}{m}\right|\le \frac{2^{A_0}}R\quad (r\text{ odd}).}
\end{equation}
In particular, for any set $D$ of odd residues modulo $2^L$,
\begin{equation}\label{eq:dangerous}
 \Psh\{n\in G:\Pass_B(n)\bmod2^L\in D\}
 \le\frac{|D|}{m}a_G+\frac{|D|2^{A_0}}R.
\end{equation}
Moreover $Q_G:=m\sum_{r\text{ odd}}q_G(r)^2$ satisfies
\begin{equation}\label{eq:goodQ}
 0\le Q_G-a_G^2\le\left(\frac{m2^{A_0}}R\right)^2.
\end{equation}
\end{lemma}
\begin{proof}
At total valuation $A$, the number of positive integer words of arbitrary length with sum $A$ equals $2^{A-1}$.  Thus there are at most $\sum_{A=1}^{A_0}2^{A-1}<2^{A_0}$ words in $G$. Apply Lemma~\ref{lem:progression} to each word and sum the discrepancies.  For \eqref{eq:goodQ}, write $q_G(r)=a_G/m+d_r$ with $\sum_r d_r=0$, and observe $m\sum_rq_G(r)^2-a_G^2=m\sum_rd_r^2\le m^2(2^{A_0}/R)^2$.
\end{proof}

\section{Controlling discarded valuation words}
Under normalized Haar measure on odd $2$-adic integers, successive Syracuse valuations $V_i$ are independent with
\[ \Ph(V_i=a)=2^{-a},\quad a\ge1,\qquad A_H=V_1+\cdots+V_H.
\]
For $0<s<1$ the geometric moment formula is
\[ \mathbb E 2^{sV_1}=\frac{2^s}{2-2^s},\qquad
\mathbb E 2^{-sV_1}=\frac{2^{-s}}{2-2^{-s}}.\]
We will use $s=1/2$ and the constants
\[
 \rho_+=\sqrt2+1,\qquad \rho_-=(2\sqrt2-1)^{-1},\qquad
 \lambda_-=-\log(\sqrt3\rho_-)=0.054149958\ldots>0.
\]

\begin{lemma}[Finite-shell and Haar comparison]\label{lem:haar}
For any integer $K\ge H$ and each event $E$ of the form $\{A_H\le K\}$ or its complement,
\[
 \left|\Psh(E)-\Ph(E)\right|\ll \frac{2^K}{B}.
\]
Consequently, if $H=\lfloor\varepsilon\log X\rfloor$, $A_0=\lfloor\gamma\log_2 X\rfloor$ and $\varepsilon\log3<\gamma$, then
\begin{align}
 \Psh(A_H>A_0)&\ll X^{-\gamma/2+\varepsilon\log\rho_+}+X^{\gamma-b},\label{eq:largeA}\\
 \Psh\left(A_H\le H\log_2 3+2\right)&\ll X^{-\varepsilon\lambda_-}+X^{\gamma-b}.\label{eq:smallA}
\end{align}
\end{lemma}
\begin{proof}
Every exact word $(a_1,\ldots,a_H)$ with total valuation at most $K$ corresponds to an odd $2$-adic cylinder of modulus $2^{A_H+1}$, hence the event $\{A_H\le K\}$ is a union of odd residue classes modulo $2^{K+1}$.  Any such set has discrepancy $O(2^K/R)=O(2^K/B)$ between normalized uniform sampling on $\mathcal I_B$ and normalized $2$-adic Haar measure.  Complements have the same bound.  Chernoff's inequality gives
\[
 \Ph(A_H>A_0)\le2^{-A_0/2}\rho_+^H,
 \qquad \Ph(A_H\le H\log_2 3+2)
 \le2^{(H\log_2 3+2)/2}\rho_-^H=2e^{-\lambda_-H}.
\]
For the second event take $K=\lfloor H\log_2 3+2\rfloor$, which is at most $A_0$ for large $X$ under the parameter hypothesis.  Floor and ceiling effects are absorbed in constants.
\end{proof}

\begin{lemma}[Long passages require atypically few divisions]\label{lem:late}
For $n\in\mathcal I_B$, if $\tau_B(n)>H$ then
\begin{equation}\label{eq:lowtail}
 A_H(n)\le H\log_2 3+2
\end{equation}
for sufficiently large $X$ and $H=O(\log X)$.
\end{lemma}
\begin{proof}
All $\Syr^j(n)>B$ for $0\le j\le H$; consequently
\[
 \Syr^H(n)=n\,\frac{3^H}{2^{A_H}}\prod_{j=0}^{H-1}
 \left(1+\frac1{3\Syr^j(n)}\right)
 \le 2B\frac{3^H}{2^{A_H}}\exp\!\left(\frac H{3B}\right).
\]
The left-hand side exceeds $B$.  Taking base-two logarithms gives
$A_H<H\log_2 3+1+H/(3B\log2)\le H\log_2 3+2$ for large $X$.
\end{proof}

\section{A local post-entrance power saving}
\begin{theorem}[Local first-passage short-excursion estimate]\label{thm:local}
Let $0<\beta< (1-\log2)/2$, let $b=1-2\beta$, and choose constants $\gamma,\varepsilon>0$ such that
\[
 \gamma<b-\log2,\qquad \varepsilon\log3<\gamma,
 \qquad \gamma/2>\varepsilon\log(\sqrt2+1).
\]
Draw $n$ uniformly from $\mathcal I_B$.  Define $\Pass_B(n)$ only when $\tau_B(n)<\infty$.  Then, with $H=\lfloor\varepsilon\log X\rfloor$,
\begin{align}
 &\Psh\Bigl(\tau_B(n)<\infty,\ 
       \max_{1\le t\le L}T^t(\Pass_B(n))>Y\Bigr)\label{eq:local}\nonumber\\
 &\qquad\ll X^{-\beta}+X^{-(b-\log2-\gamma)}
     +X^{-(\gamma/2-\varepsilon\log\rho_+)}
     +X^{-\varepsilon\lambda_-}.
\end{align}
The same bound, up to a factor of $2$, holds when the input is sampled with harmonic weight proportional to $1/n$ on $\mathcal I_B$.
\end{theorem}
\begin{proof}
For odd $y$ let $b_0,\ldots,b_{L-1}$ denote the shortcut parities under $T$ and define $Z_t(y)=3^{\sum_{j<t}b_j}/2^t$.  The odd residues modulo $2^L$ are in bijection with binary parity words beginning with $b_0=1$.  Under the uniform law on these residues, $(Z_t)_{1\le t\le L}$ is a nonnegative martingale with $Z_1=3/2$.  Also $T^t(y)=yZ_t(y)+E_t(y)$ with $0\le E_t(y)\le(3/2)^t-1$.  Since $\log(3/2)<1-\beta$, the additive error is smaller than $Y/2$ for $t\le L$ and large $X$. Thus for $y\le B$, an excursion $T^t(y)>Y$ implies $Z_t(y)>X^\beta/2$.  By Doob's inequality the set $D_X$ of dangerous odd residues has
\begin{equation}\label{eq:doob}
 \frac{|D_X|}{m}\le3X^{-\beta}.
\end{equation}
Let $G$ be the truncated, bounded-total-valuation first-passage set of Lemma~\ref{lem:balance}.  That lemma and \eqref{eq:doob} give
\[
 \Psh\{n\in G:\Pass_B(n)\bmod2^L\in D_X\}
 \le3X^{-\beta}+O\left(\frac{m2^{A_0}}B\right)
 \ll X^{-\beta}+X^{-(b-\log2-\gamma)}.
\]
Every remaining input either satisfies $A_H>A_0$ or $\tau_B(n)>H$.  Indeed, a passage by time $H$ with total valuation exceeding $A_0$ also has $A_H>A_0$.  The first exceptional event is bounded by \eqref{eq:largeA}.  By Lemma~\ref{lem:late} and \eqref{eq:smallA}, the second is $\ll X^{-\varepsilon\lambda_-}+X^{\gamma-b}$.  The $X^{\gamma-b}$ terms are absorbed by $X^{-(b-\log2-\gamma)}$.  This establishes \eqref{eq:local}.  In $(B,2B]$ the ratio between the maximum and minimum harmonic weight is at most $2$, so harmonic sampling changes any event probability by at most a factor of $2$.
\end{proof}

\begin{corollary}[Explicit power saving at the paper's parameters]\label{cor:concrete}
Set $\beta=1/200$, $\gamma=1/4$ and $\varepsilon=1/10$. Then
\[
 \boxed{\Psh\!\left(\tau_B<\infty,\ \max_{1\le t\le\lfloor\log X\rfloor}
 T^t(\Pass_B(n))>X^{199/200}\right)=O(X^{-1/200}).}
\]
Consequently the expected weighted $W_X=(X^{199/200},X]$ occupation during the first $L$ shortcut steps after passage is
\[
 \boxed{\mathbb E_{\mathrm{sh}}\sum_{t=1}^L T^t(\Pass_B(n))\mathbf1_{W_X}(T^t(\Pass_B(n)))
 =O\!\left(X^{199/200}\log X\right),}
\]
where the sum is set to zero when no first passage exists.  These estimates remain true under harmonic weighting on the shell.
\end{corollary}
\begin{proof}
Here $b=0.99$, and the positive exponents on the right of \eqref{eq:local} are
\[
\begin{aligned}
  \beta &=0.005, & b-\log2-\gamma&=0.0468528\ldots,\\
  \gamma/2-\varepsilon\log\rho_+&=0.0368626\ldots,
  &\varepsilon\lambda_-&=0.00541499\ldots.
\end{aligned}
\]
Thus the smallest exponent is $\beta=0.005$.  On the excursion event each term of the restricted window sum is at most $X$, and there are at most $L$ terms.
\end{proof}


\section{The exact source-scale obstruction and a renewal identity}
\begin{proposition}[Dyadic Ces\`aro representation of the limiting law]\label{prop:cesaro}
For fixed $B$, let $D$ be any set of odd integers not exceeding $B$, and assume that its logarithmic first-passage mass $\nu_B(D)$ exists in the sense of the source manuscript.  Write $\mathcal I_{j,B}=(2^jB,2^{j+1}B]\cap(2\N+1)$ and set
\[
 p_{j,B}(D)=\frac{\displaystyle\sum_{n\in\mathcal I_{j,B}}n^{-1}\mathbf1_{\{\tau_B(n)<\infty,\ \Pass_B(n)\in D\}}}
 {\displaystyle\sum_{n\in\mathcal I_{j,B}}n^{-1}}.
\]
Then the following identity is exact in the limiting sense:
\begin{equation}\label{eq:cesaro}
 \boxed{\nu_B(D)=\lim_{J\to\infty}\frac1J\sum_{j=0}^{J-1}p_{j,B}(D).}
\end{equation}
In particular, controlling only the initial shell $j=0$, or any fixed number of source shells, does not bound $\nu_B(D)$.
\end{proposition}
\begin{proof}
For any bounded $F$ on odd positive integers, the decomposition $w=2^in$ with $n$ odd yields
\[
 \sum_{w\le N}\frac{F(w^{\flat})}{w}
 =\sum_{\substack{n\le N\\n\ \mathrm{odd}}}\frac{F(n)}n
 \sum_{0\le i\le\lfloor\log_2(N/n)\rfloor}2^{-i}
 =2\sum_{\substack{n\le N\\n\ \mathrm{odd}}}\frac{F(n)}n+O(1).
\]
Indeed the discarded tail is at most $2n/N$ for each $n$, resulting in an $O(1)$ total error.  Apply this to the first-passage indicator defining $\nu_B(D)$, and choose $N=2^JB$. The harmonic mass of each odd dyadic shell equals $(\log2)/2+O((2^jB)^{-1})$, with summable errors.  The finitely many odd integers $n\le B$ have a fixed total harmonic mass, which disappears after division by $\log N$.  Thus the logarithmic average is the Ces\`aro mean of the shell probabilities, proving \eqref{eq:cesaro}.
\end{proof}

\begin{proposition}[First-passage renewal across one threshold]\label{prop:renewal}
Assume that the logarithmic first-passage laws $\nu_B$ and $\nu_{2B}$ exist. Then, for any set $D\subseteq[1,B]\cap(2\N+1)$,
\begin{equation}\label{eq:renewal}
 \boxed{\nu_B(D)=\nu_{2B}(D)+
 \sum_{\substack{B<z\le2B\\z\ \mathrm{odd}}}\nu_{2B}(z)
 \mathbf1_{\{\tau_B(z)<\infty,\ \Pass_B(z)\in D\}}.}
\end{equation}
\end{proposition}
\begin{proof}
Any orbit that first enters $[1,B]$ necessarily first enters $[1,2B]$.  If the latter entrance already belongs to $[1,B]$, the two entrance values agree; otherwise, the first entry $z$ is in $(B,2B]$ and the subsequent deterministic evolution to $[1,B]$ gives the second term.  The two cases partition the event pointwise. Taking the logarithmic means gives the identity because, for fixed $B$, it contains only finitely many entry states.
\end{proof}
\begin{corollary}[A concrete sufficient condition for transfer]\label{cor:transfer}
Let $D_X$ be the dangerous odd-residue set of \eqref{eq:doob}, and interpret $D_X$ as a subset of odd positive integers by reduction modulo $2^L$. Suppose for some constants $K>0$ and $\delta>0$ that
\[
  \nu_{2B}(D_X\cap[1,B])\le KX^{-\delta},
 \qquad \sup_{\substack{B<z\le2B\\z\ \mathrm{odd}}} B\,\nu_{2B}(z)\le K.
\]
Then $\nu_B(D_X\cap[1,B])\ll K(X^{-\delta}+X^{-\beta})$.
\end{corollary}
\begin{proof}
In the renewal identity \eqref{eq:renewal}, the first term is $\le KX^{-\delta}$.  In the second term, each input $z$ has weight at most $K/B$, while Theorem~\ref{thm:local} bounds the number of relevant odd $z\in(B,2B]$ by $O(BX^{-\beta})$.  Summing yields $O(KX^{-\beta})$.
\end{proof}
\begin{remark}
The pointwise incoming-weight hypothesis is a \emph{strong, unproved sufficient condition}, not a prediction from existing potential bounds.  One can replace it with a weaker estimate for incoming mass on the particular pullback of $D_X$; the difficulty remains to prove this for the actual large-source first-passage law.
\end{remark}

\begin{remark}
Equation \eqref{eq:renewal} exposes two tasks absent from the local-shell estimate: the \emph{direct overshoot} $\nu_{2B}(D)$, and the generally nonuniform incoming weights $\nu_{2B}(z)$ on $(B,2B]$.  Our local theorem controls the second term only when these incoming weights are suitably regular (for example, uniformly bounded by $C/B$).  Nothing proved here bounds the direct overshoot.  This is a more informative obstruction than assuming that the incoming states are automatically uniform.
\end{remark}

\section{What has, and has not, been established}
Let $\nu_X$ be the logarithmic limiting first-passage law from the source manuscript, obtained from starts tending to infinity before $X$ is sent to infinity.  The above proof begins with a nearly uniform source shell $(B,2B]$, and it does not say that the \emph{last pre-entry state} of a much longer orbit is so distributed.  In particular, it does not prove $Q_X=O(1)$ for $\nu_X$, does not bound the late return excess, and does not imply a polynomial bound on all Collatz orbits.

The conditional collision criterion of the preceding research note is stronger than needed to control short visits.  Here we directly bound the mass of the \emph{specific dangerous class set} $D_X$; the discarded first-passage mass can have insufficient $L^2$ regularity while still being $O(X^{-1/200})$ in total mass.  The only result about the local collision norm $Q_G$ is \eqref{eq:goodQ}; the omitted exceptional part is not controlled in that norm.

\begin{problem}[Transfer to the true entrance law]
Prove, for the actual logarithmic entrance law $\nu_X$ and the exact dangerous set $D_X$ from \eqref{eq:doob}, a one-set estimate
\[
  \boxed{\nu_X\{y:y\bmod2^L\in D_X\}=O(X^{-\delta})}
 \quad\text{for some }\delta>0.
\]
Unlike a uniform $L^2$ Fourier bound, this requests mixing only against a \emph{specific set of parity cylinders}.  A potential approach is to compare entrance laws from successive multiplicative source shells in this one-set norm, and to sum the errors; elementary progression balance gives the local ingredient but not the required inter-shell transfer or the long-time return estimate.
\end{problem}

\section{Finite data: diagnostic only}
An exhaustive C++ check for all odd starts in $(B,2B]$ for three finite values of $X$ gave the following.  Here $\widehat Q$ is the exact collision energy of the \emph{normalized finite local-shell entrance histogram}, computed directly as $m\sum_r(c_r/R)^2$ from exhaustive shell counts; it is not $Q_X$ for the limiting law.
\[
\begin{array}{c|r|r|r|r}
X & B & \Psh(\tau_B\le5) & \widehat Q & \Psh(\text{early excursion})\\\hline
2^{20} & 912{,}838 &0.75430&1.01799&0.54985\\
2^{22} &3{,}601{,}084&0.75431&1.01299&0.55244\\
2^{24} &14{,}206{,}029&0.75431&1.00591&0.55225
\end{array}
\]
The asymptotic exponent $1/200$ is extremely small, so no useful finite-$X$ numerical decay is expected at these scales.  The source code records the exact sampling and counting convention.  The local theorem is proved independently of these measurements.

\noindent\textbf{References and positioning.} The parity-word bijection and finite cylinder balance are classical (Terras--Everett; see also Lagarias' survey); the source manuscript itself gives a related finite-horizon balance lemma.  Tao's work on almost-all Collatz orbits treats a different, global first-passage stabilization problem involving $3$-adic Fourier analysis.  The local statement above repackages elementary cylinder and Chernoff estimates into a targeted bound; no priority or novelty claim is made without external review.
\end{document}
